% !TeX program = xelatex
\documentclass[11pt,a4paper]{article}
\usepackage[margin=23mm]{geometry}
\usepackage{fontspec}
\setmainfont{Latin Modern Roman}
\usepackage{amsmath,amssymb,mathtools}
\usepackage{xcolor,enumitem,booktabs,longtable,array,xurl,hyperref}
\definecolor{muted}{HTML}{53636C}
\hypersetup{colorlinks=true,urlcolor=blue,linkcolor=blue}
\setlength{\parindent}{0pt}
\setlength{\parskip}{5pt}
\newcommand{\R}{\mathbb R}
\newcommand{\E}{\mathbb E}
\newcommand{\N}{\mathbb N}
\newcommand{\ind}{\mathbf 1}
\newcommand{\Var}{\operatorname{Var}}
\newcommand{\Cov}{\operatorname{Cov}}
\newcommand{\supp}{\operatorname{supp}}
\DeclareMathOperator*{\argmin}{arg\,min}
\DeclareMathOperator*{\argmax}{arg\,max}
\newcommand{\entrymeta}[3]{{\sffamily\small\color{muted}\textbf{\texttt{#1}}\quad|\quad #2\par #3\par}}
\newcommand{\Problemref}[1]{\texttt{#1}}
\newcommand{\JELref}[2]{\texttt{#2} (JEL \texttt{#1})}
\begin{document}
\section*{Validation of panel estimators under feedback selection}
\textbf{Problem ID:} \texttt{C23-305}\quad \textbf{JEL:} \texttt{C23}\par
% BEGIN ECONOMICS STATEMENT
\label{problem:OP-0600}


\label{audited:ECON-015}

{\sffamily\small\color{muted}\textbf{ECON-015} | R | Source-backed agenda; exact formalization is editorial\par}

\subsection*{Definitions and formal target}

Observe a panel sample \(O_n\). A structural model \(m\in\mathcal M_n\) specifies treatment, potential outcomes and dependence, including untreated outcomes \(Y_{it}(0)=\alpha_i+\lambda_t+u_{it}\) with normalized errors and permitted treatment feedback from lagged \(u\). Let \(\tau(m)\) be a declared causal target, \(\widehat\tau_n(O_n)\) an estimator and \(V_n(O_n)\) a predeclared pre-treatment diagnostic. For acceptance region \(\mathcal A\), \(\varepsilon>0\) and \(\eta\in(0,1]\), study
\[
 \sup_{\substack{m\in\mathcal M_n:\\P_m(V_n\in\mathcal A)\geq\eta}}
 P_m\bigl(|\widehat\tau_n-\tau(m)|>\varepsilon\mid V_n\in\mathcal A\bigr).
\]
Characterize restrictions, validation designs or additional data yielding an informative bound, rather than treating an accepted placebo test as a certificate of causal accuracy.

\subsection*{Clarifications and required assumptions}

For each evaluation of the displayed criterion, use the model class $\mathcal M_n$ throughout, with $P_m$ the sampling law induced by $m$ and $\tau(m)\in\R$ its declared scalar causal target; writing $\tau(m)$ does not assume identification from $P_m$. The estimator and diagnostic must be measurable functions on the declared observed sample space, with the estimator finite and defined also on samples where fitting fails. Any procedure randomization must be included in the probability law. Hold these procedures and the measurable acceptance region $\mathcal A$ fixed while taking the supremum; any choice of a validation design is a separate part of the characterization, not an additional optimization implicit in the display. Conditional diagnostic error is defined for each retained model because $P_m(V_n\in\mathcal A)\geq\eta>0$. This restriction applies only to the retained subfamily and does not guarantee acceptance probability at least $\eta$ throughout $\mathcal M_n$. If that subfamily is empty, report that no models qualify; do not report zero conditional risk. At $\eta=1$, only models with almost-sure acceptance qualify. The displayed loss is the indicator of $|\widehat\tau_n-\tau(m)|>\varepsilon$, which is bounded by one regardless of outcome moments. Moment or tail assumptions may help establish useful bounds, but are not needed merely to define this probability. Expected absolute or squared error would be a different target requiring separately stated integrability conditions; bounded outcomes alone do not bound an arbitrary estimator.

An evaluated instance must declare the meaning of $n$, panel dimensions and observation scheme, structural and observable probability spaces, sampling law, treatment histories and interventions, and the target's population, horizon and weights. It must also declare the exact error-centering condition, fixed-effect location normalizations, initial-history law, assignment information and dependence restrictions. A zero-mean normalization is not a substitute for a dependence law or an exogeneity assumption. The permitted observable information and timing for the diagnostic must be specified, especially when treatment dates vary or depend on past shocks. These choices, and the quantitative meaning of an informative bound, remain unspecified here; no particular centering condition, feedback model, diagnostic window or finite-sample or asymptotic success criterion is selected.

\subsection*{Variants and changes of scope}

\textbf{Calibration.} The displayed conditional-risk criterion is the calibration target; this label introduces no distinct loss or optimization. \textbf{External validation.} Compare with a randomized benchmark population having a declared transport relationship to the observational sample. \textbf{Feedback.} Permit selection based on lagged shocks with an explicit dependence restriction; contrast strict exogeneity. \textbf{Coverage.} Confidence-set coverage after diagnostic selection is a separate variant. A formal instance must specify the confidence level and whether the guarantee is conditional on $V_n\in\mathcal A$ or simultaneous over a declared index set. Simultaneous coverage alone does not imply the same coverage level conditional on acceptance; an ordinary marginal confidence interval likewise does not by itself establish that conditional guarantee.

\subsection*{Audit conclusion and source relationship}

Section 10.2 expressly discusses validation and limitations of placebo tests under feedback selection. The worst-case conditional-error criterion is an extension proposed by this catalogue, not a published impossibility or existence claim.

\subsection*{References and exact source roles}

{\footnotesize\raggedright
\par\noindent\textbf{[S1]} Dmitry Arkhangelsky and Guido Imbens (2024). \emph{Causal Models for Longitudinal and Panel Data: A Survey}. The Econometrics Journal 27(3), C1--C61, DOI: 10.1093/ectj/utae014. Checked author version: arXiv:2311.15458v3, 25 June 2024.\par \textit{Verified locator / source role:} Section 10.2 in the checked June 2024 author version. The journal landing page was verified for metadata; the readable author PDF was used for the agenda.\par \url{https://arxiv.org/pdf/2311.15458}\par\smallskip
}

\subsection*{Common conventions: Additional-source status and mathematical conventions}

\label{audited:conventions}

Of the 100 supplied ECON entries, 65 distinct problems appear here; 32 close
matches refine earlier entries, and three duplicate questions map to existing
entries. Appendix~\ref{app:audited-crosswalk} lists every destination.
The attachment's P/R/S/U distinctions and source audits are retained. Its
precise mathematical formalizations are editorial unless the individual audit
says otherwise. Candidate subsidy targets ECON-004--006 retain their U flags;
the resolved approval-core question ECON-007 updates OP-0202 above.
The following supplied conventions govern both this part and the attributed
ECON refinements elsewhere in the catalogue, subject to local restrictions.

\textbf{Parameterized questions.} A problem family lists inputs that must be fixed before an instance is evaluated: population, horizon, policies, information, data law, model class and welfare weights. These are required choices, not quantities the citations are assumed to specify. Undefined applications should not be treated as identified estimands.

\textbf{Indices and integrability.} Sets of agents, firms and regions are finite unless a population measure is explicitly used. \(i,j\) index units, \(r\) regions, \(t\) dates and \(h\) horizons. \(\mathbb E\) and \(\Pr\) refer to the declared population and law; \(\mathbf1\{A\}\) is an indicator. Every displayed expectation and discounted sum needs existence in the stated domain. Infinite horizons use a discount in \((0,1)\) and a convergence condition; finite horizons may use discount one.

\textbf{Optimization and existence.} A displayed \(\max\), \(\min\), or \(\arg\max\) is conditional on attainment. For finite feasible menus this is immediate; general spaces need continuity and compactness/coercivity or another existence argument. Without attainment, the target is the supremum/infimum and arbitrarily near-optimal feasible choices. \(\inf\varnothing=+\infty\). Empty feasibility and an unidentified target are different issues.

\textbf{Potential outcomes.} \(Y_i(z)\) is the outcome under a specified intervention, not the observed conditional mean among people choosing \(z\). In binary comparisons \(\Delta Y_i=Y_i(1)-Y_i(0)\). Specify assignment, treatment versions, compliance, information and observation horizon. Interference is allowed under a full policy vector; compressed exposure notation needs a justified mapping. No interference, consistency and overlap are substantive assumptions, not automatic consequences of notation.

\textbf{Populations and observation.} Fix baseline eligibility and population weights. Conditioning on post-policy employment, survival, relocation or program uptake can change the population being compared. Death is not ordinary loss to follow-up; outcomes truncated by death need a composite convention, joint survival target, or explicitly defined survivor stratum. Fertility-changing interventions need a population functional rather than an unexplained fixed descendant cohort. Measurement and missingness models must be supplied.

\textbf{Identification.} A structural model \(m\in\mathcal M\) implies observable law \(P_m\) and target \(\theta(m)\). Identification means
\[
 P_{m_1}=P_{m_2}\ \Longrightarrow\ \theta(m_1)=\theta(m_2)
 \quad\text{for all }m_1,m_2\in\mathcal M .
\]
Without identification, the target is
\[
 \Theta(P;\mathcal M)=\{\theta(m):m\in\mathcal M,\ P_m=P\}.
\]
Writing \(\theta(P)\) before establishing identification can hide incompatible latent models. Confidence sets additionally need a sampling law, confidence level, asymptotic sequence and uniformity class. Point identification, coverage and informative precision are separate properties.

\textbf{Mediation.} Total effects of randomized assignment are distinct from cross-world natural indirect effects. Belief/effort-channel effects require a meaningful mediator intervention and additional measurement, exclusion or mediation assumptions. Randomization of the initial policy alone does not supply those assumptions. Report bounds or interventional alternatives when the stronger target is unsupported.

\textbf{Equilibrium.} A counterfactual \(W(z)\), \(p(z)\), or \(\mu_t^z\) requires individual optimization, feasibility, government/resource budgets, market clearing and state-transition laws. State equilibrium existence and selection, or report ranges over compatible equilibria. Initial-state integration includes subsequent endogenous transitions; current-state integration must use the policy-dependent distribution. Reduced-form invariance after policy is not assumed.

\textbf{Welfare accounts.} Scores, utility, health, dollars and ecological quantities require explicit conversion before addition. Fees, taxes, subsidies, wages and extortion payments are transfers between parties; distributional weights can make them welfare relevant, but their face value is not automatically a resource loss. State ownership, geographic boundaries, foreign-income weights and fiscal finance. Do not add profits to household welfare already including dividends, or count land capitalization and the underlying benefit twice. A benefit-cost expression must distinguish gross benefits from net welfare and subtract every real cost exactly once.

\textbf{Derivatives and risk.} Log derivatives require positive arguments; local responses require admissible perturbations and specified equilibrium branches. Differentiation under expectations needs regularity/domination, and boundaries may allow only one-sided derivatives. Risk uses a specified law; ambiguity uses a declared nonempty law set on a common history space. Adaptive policies must be nonanticipative. A robust criterion is an editorial choice unless attributed explicitly.

\textbf{Scope overlaps.} Allocation, collective choice, contracts, household bargaining and coordinated policy overlap economics with optimization and game theory. They are retained because they appeared in the original catalogue; no claim of a strictly game-theory-free selection is made.
% END ECONOMICS STATEMENT
\end{document}
